% Appendix A -- full tables of the post-tuning re-score (registered experiment).
\chapter{Post-tuning Re-score of PFedRec: Full Tables}
\label{app:posttune}

The experiment is described in Section~\ref{sec:method-posttune} and its pre-registration is the
file \path{posttune/README.md} in the directory of run \RunPfedrecShort{} under
\path{results/federated/pfedrec/} (version 1.1, 4 November 2026); the filled result file is
\path{RESULT.md} in the same directory. Inputs: the restored
best-round item table (SHA-256 \PtGlobalsSha), the \DsUsers{} calibrated heads, and the
candidate lists of the run's own evaluation. The run took \PtElapsedSec{} s on a CPU for the
four cells. Trust checks: the replay of the reference matched the run bit-exactly; no
non-finite value in any $\Qu$; untouched rows bit-identical to $\Qtab$ on a sample of users;
the test item was drawn at least once for \PtNHitClean{} users in the clean pool and for
\PtNHitLeaky{} in the official-code pool; and no user's score decreased in the official-code
cells (the registered monotonicity check). The median row shift was \PtRowShiftMedianEOne{} at
$E=1$ and \PtRowShiftMedianEThree{} at $E=3$ against a median row norm of \PtRowNormMedian{} in
$\Qtab$.

\section{All Four Cells}

\begin{table}[htbp]
  \caption{\ndcg{} / \hr{} / \mrr{} of the reference and the four cells, overall and by group.}
  \label{tab:posttune-cells}
  \centering\scriptsize
  \fittable{% AUTO-GENERATED by thesis/scripts/build_assets.py -- do not edit by hand.
% posttune/reeval_qu_diag_best.json: one row per cell and metric, columns = overall and groups
\begin{tabular}{llrrrr}
\toprule
Cell & Metric & Overall & Sparse & Medium & Dense \\
\midrule
Shared item table (reference) & NDCG@10 & 0.3447 & 0.4270 & 0.3802 & 0.2935 \\
 & HR@10 & 0.6111 & 0.7058 & 0.6688 & 0.5387 \\
 & MRR & 0.2821 & 0.3552 & 0.3083 & 0.2409 \\
\midrule
$E=1$, clean & NDCG@10 & 0.3181 & 0.4165 & 0.3611 & 0.2564 \\
 & HR@10 & 0.5584 & 0.6848 & 0.6326 & 0.4641 \\
 & MRR & 0.2608 & 0.3474 & 0.2937 & 0.2105 \\
\midrule
$E=1$, official-code & NDCG@10 & 0.3789 & 0.4305 & 0.3891 & 0.3564 \\
 & HR@10 & 0.6642 & 0.7108 & 0.6792 & 0.6394 \\
 & MRR & 0.3090 & 0.3583 & 0.3165 & 0.2893 \\
\midrule
$E=3$, clean & NDCG@10 & 0.2762 & 0.3981 & 0.3347 & 0.1956 \\
 & HR@10 & 0.4704 & 0.6539 & 0.5792 & 0.3324 \\
 & MRR & 0.2305 & 0.3329 & 0.2742 & 0.1672 \\
\midrule
$E=3$, official-code & NDCG@10 & 0.4520 & 0.4385 & 0.4128 & 0.4871 \\
 & HR@10 & 0.7505 & 0.7219 & 0.7063 & 0.7937 \\
 & MRR & 0.3738 & 0.3649 & 0.3381 & 0.4047 \\
\bottomrule
\end{tabular}
}
\end{table}

\section{Per-group Paired Tests}

\begin{table}[htbp]
  \caption{Paired tests against the shared-table reference by user group. Where $n_-=0$ the
  one-sided $p$ follows from the construction; for groups with few changed users the computed
  value is printed.}
  \label{tab:posttune-groups}
  \centering\scriptsize
  \fittable{% AUTO-GENERATED by thesis/scripts/build_assets.py -- do not edit by hand.
% stats_wilcoxon.json cells.<cell>.by_scope.<group>; structural label when W^-=0 and >60 users changed
\begin{tabular}{llrrrrrr}
\toprule
Cell & Group & $\Delta$NDCG@10 & 95\% CI & $\Delta$HR@10 & $p$ one-sided & $p$ Pratt & $n_+/n_-/n_0$ \\
\midrule
$E=1$, clean & sparse & \mbox{-0.0105} & \mbox{[-0.0181, -0.0039]} & \mbox{-0.0210} & 0.955 & 0.124 & 30/21/758 \\
$E=1$, clean & medium & \mbox{-0.0191} & \mbox{[-0.0250, -0.0133]} & \mbox{-0.0362} & 0.64 & $<10^{-10}$ & 246/108/1968 \\
$E=1$, clean & dense & \mbox{-0.0371} & \mbox{[-0.0460, -0.0286]} & \mbox{-0.0746} & 0.999 & $<10^{-10}$ & 703/387/1819 \\
\midrule
$E=1$, official-code & sparse & +0.0036 & \mbox{[+0.0018, +0.0058]} & +0.0049 & $2.6\times10^{-7}$ & $4.6\times10^{-9}$ & 33/0/776 \\
$E=1$, official-code & medium & +0.0090 & \mbox{[+0.0073, +0.0108]} & +0.0103 & W$^-$=0 & W$^-$=0 & 262/0/2060 \\
$E=1$, official-code & dense & +0.0629 & \mbox{[+0.0583, +0.0675]} & +0.1007 & W$^-$=0 & W$^-$=0 & 1034/0/1875 \\
\midrule
$E=3$, clean & sparse & \mbox{-0.0289} & \mbox{[-0.0414, -0.0176]} & \mbox{-0.0519} & 0.971 & 0.00418 & 91/53/665 \\
$E=3$, clean & medium & \mbox{-0.0454} & \mbox{[-0.0550, -0.0363]} & \mbox{-0.0896} & 1.0 & $9.3\times10^{-9}$ & 507/279/1536 \\
$E=3$, clean & dense & \mbox{-0.0979} & \mbox{[-0.1099, -0.0858]} & \mbox{-0.2063} & 1.0 & 1.0 & 683/828/1398 \\
\midrule
$E=3$, official-code & sparse & +0.0115 & \mbox{[+0.0081, +0.0152]} & +0.0161 & W$^-$=0 & W$^-$=0 & 98/0/711 \\
$E=3$, official-code & medium & +0.0326 & \mbox{[+0.0294, +0.0361]} & +0.0375 & W$^-$=0 & W$^-$=0 & 629/0/1693 \\
$E=3$, official-code & dense & +0.1936 & \mbox{[+0.1850, +0.2020]} & +0.2551 & W$^-$=0 & W$^-$=0 & 1775/0/1134 \\
\bottomrule
\end{tabular}
}
\end{table}

\section{Registered Predictions and Outcome}

\begin{table}[htbp]
  \caption{The five predictions registered before the run and their outcome. A1 was judged on
  the mean with its confidence interval; under Pratt's treatment of ties the rank test points
  the other way (Section~\ref{sec:res-pfedrec}).}
  \label{tab:predictions}
  \centering\small
  \fittable{% AUTO-GENERATED by thesis/scripts/build_assets.py -- do not edit by hand.
% posttune/README.md section 8 predictions vs RESULT.md section 4 (numbers via macros)
\begin{tabular}{@{}p{0.04\linewidth}>{\raggedright\arraybackslash}p{0.40\linewidth}>{\raggedright\arraybackslash}p{0.36\linewidth}>{\raggedright\arraybackslash}p{0.13\linewidth}@{}}
\toprule
\# & Prediction (registered before the run) & Observed & Verdict \\
\midrule
A1 & Clean negatives, $E=1$: NDCG@10 below the reference; $E=3$ below $E=1$ & \NdcgEOneClean{} ($\Delta$ \PtDeltaEOneClean{}, CI \PtCiEOneClean{}); $E=3$: \NdcgEThreeClean{} & Correct on the mean \\
A2 & Official-code negatives, $E=1$: $\Delta \geq +0.010$, one-sided $p<0.05$; $E=3 \geq E=1$ & $\Delta$ \PtDeltaEOneLeaky{} \PtCiEOneLeaky{}; $E=3$: \NdcgEThreeLeaky{} & Correct \\
A3 & Official $-$ clean $\geq +0.010$ at both $E$ & \PtContrastEOne{} and \PtContrastEThree{} & Correct \\
A4 & Dense users have the largest $|\Delta|$ in all four cells & holds in 4/4 cells (Table~\ref{tab:posttune-groups}) & Correct \\
A5 & No cell reaches 0.394 (half of the gap to the published 0.4436) & $E=3$, official-code negatives: \NdcgEThreeLeaky{} & Refuted \\
\bottomrule
\end{tabular}
}
\end{table}

\section{Reporting Rule}

The README fixed, before the run, that the thesis table would contain three rows (reference;
clean $E=1$; official-code $E=1$) with no selection of the highest cell, that the $E=3$ cells
would appear only in this appendix, and that the limitation sentence would be chosen by the
following rule: if official $-$ clean at $E=1$ is at least $+0.010$ with one-sided $p<0.05$,
then ``part of the gap to the published number comes from the official code excluding the test
item from the negatives during post-tuning''; otherwise ``the difference in the negative pool is
negligible and the residual gap is not separated between $C$ and K1--K3''. The first branch
applied (\PtContrastEOne, CI \PtContrastEOneCi).
